\section{Specialization of the coherent Hecke system}
\label{sec:specialization}

We now prove that geometric nearby cycles preserve the entire unramified
Hecke system on the smooth leg locus.  The proof must include collisions:
commutation with a separate one-point transform would not supply the
fusion identities required in Theorem~\ref{thm:main}.  The essential local
calculation is an exterior-product comparison on a space that may already
carry other modifications and leg parameters.  The corresponding
unlevelled comparison and its fusion diagrams appear in
\cite[Proposition~4.2.5 and \S4.2.6]{GR}.  We give the local proof that
also applies with the disjoint level data used here.

Throughout this section $S,\bar\eta,s$ have the meaning fixed in
Section~\ref{sec:foundations}.  Let $\cB/S$ be a smooth bundle stack and
$U/S$ a smooth scheme of relative dimension one parametrizing allowed
legs.  We assume the usual spherical local description at those legs:
after choosing a coordinate and a frame of the input bundle, a bounded
one-step Hecke correspondence is a split Schubert model in the curve
variable.  Its maps to input-and-leg and to output-and-leg are
representable and proper on bounds.  Its open exact-position maps are
smooth.  These conditions hold for ordinary or enhanced level away from
the marked point, and for the twisted-curve family used in
Section~\ref{sec:construction}.  We verify the last assertion at the end
of the section.

\begin{proposition}\label{prop:hecke-specialization}
For $F\in D_{\mathrm{lcc}}(\cB_{\bar\eta},E)$, any finite set $I$, and
any representations $\boldsymbol V=(V_i)_{i\in I}$, there are natural
isomorphisms
\begin{equation}\label{eq:spec-hecke-comparison}
 c_{I,\boldsymbol V}(F):
 \Hecke^s_{I,\boldsymbol V}(\Psi_{\cB}F)
       \xrightarrow{\ \sim\ }
 \Psi_{\cB\times_S U^I}
                      \bigl(\Hecke^{\bar\eta}_{I,\boldsymbol V}(F)\bigr).
\end{equation}
They are compatible with the tensor unit, successive Hecke
modifications and their convolution maps, permutations, and ordinary
pullback along every collision diagonal.  The functors have the relative
normalization of \eqref{eq:found-hecke-definition}.

Suppose, moreover, that $F$ has a coherent eigenstructure with lisse tensor
eigenvalue $V\mapsto L_{\bar\eta,V}$ on $U_{\bar\eta}$.  Suppose the
lisse lattices specialize to a tensor system $V\mapsto L_{s,V}$ on $U_s$
as in Proposition~\ref{prop:nearby-foundations}.  Then $\Psi_{\cB}F$ has
a coherent eigenstructure with this special-fiber eigenvalue:
\begin{equation}\label{eq:spec-eigenvalue}
 \Hecke^s_{I,\boldsymbol V}(\Psi F)
       \simeq (\Psi F)\boxtimes
                         \boxtimes_{i\in I}L_{s,V_i}.
\end{equation}
This conclusion does not assert that $\Psi F$ is nonzero; nonvanishing
is supplied separately by Theorem~\ref{thm:specialization-detection}.
\end{proposition}

We first establish the exterior-product calculation with arbitrary input,
and then apply it to fixed-point Satake kernels and their fusion products.

\begin{lemma}[The frame comparison with arbitrary input]
\label{lem:spec-frame}
Let $\cH\to\cB\times_S U$ be a bounded one-step Hecke correspondence,
viewed by its input-and-leg map, and let $K_{\bar\eta}$ be its relative
Satake kernel.  Let $b:B'\to\cB\times_S U$ be any finite-type map of
models and set
\[
 W=B'\mathbin{\times}_{\cB\times_S U}\cH,
 \qquad r:W\to B',\qquad h:W\to\cH.
\]
For $D\in D^b_c(B'_{\bar\eta},E)$, the natural composite
\begin{align}
 \Theta_b(D,K):\quad
 r_s^*\Psi_{B'}D\otimes h_s^*\Psi_{\cH}K_{\bar\eta}
 &\xrightarrow{\operatorname{Ex}^*_r\otimes
                             \operatorname{Ex}^*_h}
 \Psi_W(r_{\bar\eta}^*D)\otimes
                      \Psi_W(h_{\bar\eta}^*K_{\bar\eta})\notag\\
 &\xrightarrow{\mu}
 \Psi_W(r_{\bar\eta}^*D\otimes
                            h_{\bar\eta}^*K_{\bar\eta})
                       \label{eq:spec-frame-map}
\end{align}
is an isomorphism.  The analogous statement for twisted products is
obtained by frame descent.  No smoothness assumption on $B'$ or on $b$
is required.
\end{lemma}

\begin{proof}
All assertions are local for smooth covers.  Choose a smooth cover
$C\to\cB\times_S U$ carrying a coordinate at the leg and an input
frame to a sufficiently high jet order for the bound.  Since
$\cB\times_S U$ is smooth over $S$, we may take $C$ smooth over $S$.
In these coordinates
\[
   \cH_C\simeq C\times_S\Gr_{\leq\lambda,S},\qquad
   W_C\simeq B'_C\times_S\Gr_{\leq\lambda,S},
 \quad B'_C=B'\mathbin{\times}_{\cB\times_S U}C,
\]
enlarging the bound if the kernel is a sum of simple Satake objects.
The pullback of $K_{\bar\eta}$ is the pullback of the Grassmannian
Satake object $K^{\Gr}_{\bar\eta}$.  As $C/S$ is smooth, smooth
exchange identifies its nearby cycles with the pullback of
$\Psi_{\Gr}K^{\Gr}_{\bar\eta}$.  Smooth exchange also identifies
the pullback of $\Psi_{B'}D$ to $B'_{C,s}$ with the nearby cycles of
the pullback $D_C$.

Under these identifications, \eqref{eq:spec-frame-map} is exactly
\[
 \Psi_{B'_C}D_C\boxtimes\Psi_{\Gr}K^{\Gr}_{\bar\eta}
     \longrightarrow
 \Psi_{B'_C\times_S\Gr}
                       (D_C\boxtimes K^{\Gr}_{\bar\eta}).
\]
It is an isomorphism by the exterior-product theorem in
Proposition~\ref{prop:nearby-foundations}, which permits a nonsmooth
first factor.  The calculation is compatible with changes of frame
because every arrow is a natural exchange map.  It therefore descends.
Notice that the two individual pull-exchange arrows in
\eqref{eq:spec-frame-map} were not required to be isomorphisms; the
assertion concerns their composite with tensor exchange.
\end{proof}

\begin{lemma}\label{lem:spec-satake}
Let $\Gr_{\leq\lambda,S}$ be the split projective Schubert model for
loops in the curve variable over $S$.  Then
\begin{equation}\label{eq:spec-simple-satake}
                 \Psi\IC_{\lambda,\bar\eta}
                                  \simeq\IC_{\lambda,s}.
\end{equation}
Specialization is a symmetric tensor equivalence of the two spherical
Satake categories and intertwines their cohomological fiber functors.
Under Satake equivalence, it is tensor isomorphic to the identity on
$\Rep(\Gd)$.  Fix such a tensor identification in the rest of the paper.
\end{lemma}

\begin{proof}
Let $P=\Psi\IC_{\lambda,\bar\eta}$.  By
Proposition~\ref{prop:nearby-foundations}, $P$ is perverse.  It is
equivariant under the positive-loop group: on a bound the action and its
equivariance diagram are computed at a sufficiently large finite jet
level, where the relevant projections are smooth, and their equivariance
maps specialize by smooth pullback.  The same observation applies to
coordinate changes.  The open Schubert orbit is smooth over $S$, so
\[
               P|_{\Gr_s^\lambda}=E[d_\lambda](d_\lambda/2).
\]
Its support is contained in $\Gr_{\leq\lambda,s}$.

The characteristic-zero coefficient Satake category on either geometric
fiber is semisimple and is identified with $\Rep(\Gd)$, with simples
indexed by the same dominant coweights.  This is geometric Satake over
an arbitrary algebraically closed base field
\cite[Theorem~14.1]{MV}.  Thus $\IC_{\lambda,s}$ occurs in $P$ with
multiplicity one.  Any other constituent has smaller highest weight.
Proper nearby-cycle compatibility gives
\begin{equation}\label{eq:spec-satake-cohomology}
 R\Gamma(\Gr_{\leq\lambda,s},P)
       \simeq
 R\Gamma(\Gr_{\leq\lambda,\bar\eta},
                                      \IC_{\lambda,\bar\eta}).
\end{equation}
The total dimension on the right is $\dim V_\lambda$, where
$V_\lambda$ is the irreducible $\Gd$-representation with highest weight
$\lambda$.  The constituent $\IC_{\lambda,s}$ already has that total
cohomological dimension.  Every nonzero Satake constituent has nonzero
cohomology, since total cohomology is its faithful exact fiber functor
\cite[\S6]{MV}.  There can be no additional constituent, proving
\eqref{eq:spec-simple-satake} with the stated normalization.

For convolution use the usual proper map from the twisted product of
two Schubert models.  Lemma~\ref{lem:spec-frame}, with the first
modification as its input space, identifies nearby cycles of the twisted
product of kernels with the twisted product of their nearby cycles.
Proper exchange then gives
\begin{equation}\label{eq:spec-convolution-kernels}
 \Psi K\star\Psi L\xrightarrow{\ \sim\ }\Psi(K\star L).
\end{equation}
Associativity and the unit constraint follow from the compatibility of
the pull, tensor, and proper-push exchange maps.  Apply the frame lemma
also to the global successive-modification diagram over $U^2$ and then
push by its proper convolution map.  This first identifies the
specialization of the entire fused kernel with the special-fiber fused
kernel, including the diagonal.  It then compares the fusion
commutativity constraints: off the diagonal the map is exterior symmetry;
on the local Grassmannian model, after adding the two leg shifts, the
resulting special-fiber fusion kernels are the perverse middle extensions
from that open locus.  Any additional smooth chart factor carries its
usual smooth-pullback shift.  Their
maps are therefore determined there.  No commutation of nearby cycles
with arbitrary middle extension is used.  This is the construction of
the symmetry in \cite[\S5, especially~(5.10)--(5.11)]{MV}, and the same
argument with more legs supplies its coherence.  The conventional parity
adjustment in \cite[\S6]{MV} is trivial here, since all Schubert
dimensions are even.

Consequently $\Psi$ gives a symmetric tensor equivalence of the two
Satake categories, preserving their simple labels and their
cohomological fiber functors by
\eqref{eq:spec-satake-cohomology}.  Under the fixed Satake equivalences,
it is a symmetric tensor autoequivalence of $\Rep(\Gd)$ preserving every
highest weight.  Tannakian duality identifies it, up to tensor
isomorphism, with an automorphism of $\Gd$ up to inner automorphism.
Preservation of every highest weight makes the induced automorphism of
the based root datum trivial, so this automorphism is inner.  Over the
algebraically closed coefficient field this supplies a tensor
isomorphism to the identity.  We choose it once; all further comparisons
use that choice rather than separate identifications for individual
representations.  This choice need not induce the previously specified
cohomology comparison in a fixed trivialization of the fiber functor.
The possible change is inner conjugacy, which does not change the
isomorphism class of a $\Gd$-local system.
\end{proof}

\begin{proof}[Proof of Proposition~\ref{prop:hecke-specialization}]
For one leg, first pull $F$ to $\cB_{\bar\eta}\times U_{\bar\eta}$.
The projection to $\cB$ is smooth, so its pull exchange is an
isomorphism.  Apply Lemma~\ref{lem:spec-frame} with
$B'=\cB\times_S U$, and identify the specialized kernel by
Lemma~\ref{lem:spec-satake}.  Proper push exchange along the
output-and-leg map gives the comparison in the orientation
\[
 \Hecke^s_V(\Psi F)
   \longrightarrow o_{s,*}\Psi_{\cH}
                     (F\,\widetilde\boxtimes\,K_{\bar\eta,V})
   \xrightarrow{\operatorname{Ex}_{o,*}^{-1}}
                     \Psi_{\cB\times_SU}\Hecke^{\bar\eta}_V(F).
\]
Both arrows are isomorphisms.  All constructions are made on a bound
containing the kernel's support and are unchanged upon increasing it.

For several legs, choose an order $i_1,\ldots,i_m$ of $I$ and form the
stack of successive modifications.  At step $r$ it remembers bundles
$P_0,\ldots,P_r$, the leg tuple, and the first $r$ modifications.
The next correspondence is attached by the bundle $P_r$ and the next
leg.  The space $B'$ in Lemma~\ref{lem:spec-frame} is now this space
of preceding data.  Its complex is the previous twisted product, and
it can be singular.  The lemma proves inductively that the product of
the pulls of $\Psi F$ and of the specialized step kernels maps
isomorphically to nearby cycles of the successive twisted product.
This remains true if some or all of the legs have been set equal:
the lemma permits any map $B'\to\cB\times_S U$.

Forget the intermediate bundles.  This convolution map is proper on the
chosen bounds.  Proper exchange and the projection formula compare its
pushforward with the transform defined by the fused kernel
$\operatorname{Sat}_I(\boldsymbol V)$.  This produces
\eqref{eq:spec-hecke-comparison}.  Alternatively one can push each
intermediate output before performing the next modification.  The two
comparisons agree: pull-and-tensor exchange is associative, proper-push
exchange composes, and the proper base-change and projection-formula
identities move an intermediate push through the next twisted product.
Thus the result is compatible with the succession of Hecke functors,
not merely with the final objects of the two constructions.

We give the diagonal comparison explicitly.  Fix
$a:I\twoheadrightarrow J$ and write $\delta=1\times\delta_a$.
For a sheaf $Q$ on $\cB_{\bar\eta}\times U_{\bar\eta}^I$, the
natural diagonal exchange has the form
\begin{equation}\label{eq:spec-diagonal-exchange}
        \delta_s^*\Psi Q\longrightarrow\Psi(\delta_{\bar\eta}^*Q).
\end{equation}
It need not be an isomorphism for arbitrary $Q$.  We prove that it is
an isomorphism for the complexes in the Hecke construction, and that it
identifies their comparison maps.

Before the first modification, $Q$ is the pullback of $F$ from $\cB$.
Both projections $\cB\times_S U^I\to\cB$ and
$\cB\times_S U^J\to\cB$ are smooth.  Compatibility of pull exchange
for their factorization through $\delta$ identifies
\eqref{eq:spec-diagonal-exchange} with the isomorphism between the two
smooth-pullback formulas.  Suppose the assertion has been proved on the
space of the first $r-1$ modifications.  Pull a coordinate-and-frame
cover for the next step to the diagonal.  The next twisted product on
each side is, in these covers, the exterior product of the preceding
complex with the fixed Grassmannian Satake kernel.  The square relating
the two maps \eqref{eq:spec-frame-map} and diagonal exchange commutes
by naturality.  The map for the first factor is an isomorphism by the
induction hypothesis, and both frame maps are isomorphisms by
Lemma~\ref{lem:spec-frame}.  The diagonal map for the new twisted
product is therefore an isomorphism too.  Proper base change and proper
exchange preserve this conclusion when intermediate bundles are
forgotten.  Induction proves it for the full transform and proves that
the comparison agrees with the convolution identification on the
collision diagonal.  Applied to the kernel construction itself, this
also proves the assertion for the fused Satake kernel used in
\eqref{eq:found-hecke-definition}.

For clarity, the square just established at the level of Hecke transforms
is
\[
\begin{CD}
 \delta_s^*\Hecke^s_{I,\boldsymbol V}(\Psi F)
       @>{\delta_s^*c_{I,\boldsymbol V}}>>
             \delta_s^*\Psi \Hecke^{\bar\eta}_{I,\boldsymbol V}(F)\\
 @V{\mathrm{fusion}}V{\sim}V
       @VV{\operatorname{Ex}^*_{\delta}\,,\ \Psi(\mathrm{fusion})}V\\
 \Hecke^s_{J,\boldsymbol W}(\Psi F)
       @>{c_{J,\boldsymbol W}}>>
             \Psi \Hecke^{\bar\eta}_{J,\boldsymbol W}(F),
\end{CD}
\]
where $W_j=\bigotimes_{i\in a^{-1}(j)}V_i$.  The right vertical
arrow is an isomorphism by the preceding induction.  For two successive
surjections, the square for their composite is the composite of the two
squares, because pull exchange composes.  This checks nested collisions
as well as a single diagonal.

It remains to identify permutations and the unit.  Over the locus of
pairwise distinct legs, the comparison is an exterior-product
comparison, and hence respects every permutation.  The permutation
maps on the relative fused Satake kernels are determined from this
locus by middle extension on the Grassmannian models after the leg shift.
The kernel comparison
therefore respects them everywhere, as in
Lemma~\ref{lem:spec-satake}.  Applying the same twisted product with
$F$ and the same proper push carries this equality of kernel maps to
the Hecke transform; no assertion that the transform of $F$ itself is
a middle extension is needed.  This gives independence of the order
chosen for successive modifications and all symmetric-group
relations.  For the trivial representation the kernel is supported on
the identity modification, so the comparison is simply smooth exchange
for $\cB\times_S U^I\to\cB$.  Its unit constraint is the unit
constraint of the exchange maps.  Naturality in every representation
has held throughout, because every construction was performed on the
Satake category and on maps of kernels.

Finally apply $\Psi$ to the given eigenisomorphism and compose with
\eqref{eq:spec-hecke-comparison}.  Smooth pullback in the leg variables
and \eqref{eq:found-lisse-tensor} identify its target with the
right-hand side of \eqref{eq:spec-eigenvalue}.  On a collision diagonal,
the same lisse tensor formula identifies the product of eigenvalue
factors with $L_{s,\otimes V_i}$.  Hence diagonal exchange is an
isomorphism on the eigenvalue side as well.  Apply these comparisons
to each original coherence diagram.  The comparisons commute with its
arrows by the preceding construction, so its specialized diagram also
commutes.  This proves unit, tensor, permutation, and fusion coherence
of \eqref{eq:spec-eigenvalue}.
\end{proof}

\subsection{Why disjoint level data do not alter the comparison}

We spell out the locality assertion used in the proposition.  A
modification at a leg $u\in U$ is obtained by gluing a bundle on the
complement of its graph to a bundle on the formal disc at $u$, with an
identification on the punctured disc.  The level structure at a disjoint
marked point belongs entirely to the first piece and is transported by
that identification.  A frame on the formal disc therefore gives the
same Schubert model with ordinary, enhanced, or no level structure.
The frame and coordinate choices needed on a fixed bound reduce to
finite jet torsors, which are smooth.  From the output side, inversion
of the modification gives an opposite Schubert bound; hence the
output-and-leg map is proper as well.  These are the usual
Beauville--Laszlo descriptions of the Hecke correspondence.

For a twisted nodal curve, choose the leg in its smooth scheme locus.
The same gluing takes place on an ordinary formal disc and leaves the
stacky node in the complementary piece.  In particular it preserves
the isomorphism class of the stabilizer representation at the node.
Restricting the bundle stack by deleting unwanted closed-fiber types
therefore restricts the Hecke correspondence by the same condition on
either side; its bounded maps remain proper by base change.  Coordinates,
frames, and the open exact-position smoothness are unchanged at the
leg.  After the nodal identification
\eqref{eq:found-node-quotient}, the action on either punctured component
is the action on that enhanced-flag factor, followed by the simultaneous
torus quotient.  This establishes precisely the local hypotheses of
Proposition~\ref{prop:hecke-specialization} in both specializations
used below.
